% !TEX program = xelatex
\documentclass[12pt,a4paper]{article}

% ========== بسته‌های پایه ==========
\usepackage{amsmath,amssymb,amsthm}
\usepackage{geometry}
\usepackage{enumitem}
\usepackage{xcolor}
\usepackage{graphicx}
\usepackage{setspace}
\usepackage{fancyhdr}
\usepackage{titlesec}

% ========== فونت فارسی ==========
\usepackage{fontspec}
\usepackage{xepersian}

\settextfont[Scale=1.05]{XB Niloofar}
\setlatintextfont[Scale=1]{Times New Roman}
\setdigitfont[Scale=1.05]{XB Niloofar}

% ========== تنظیمات صفحه ==========
\geometry{margin=2.5cm, top=3cm, bottom=2.5cm}
\setlength{\parindent}{0pt}
\setlength{\parskip}{0.8em}
\onehalfspacing

% ========== رنگ‌ها ==========
\definecolor{mainblue}{RGB}{25, 55, 120}
\definecolor{darkgreen}{RGB}{20, 100, 60}
\definecolor{darkred}{RGB}{140, 30, 30}

% ========== سبک عنوان‌ها ==========
\titleformat{\section}
  {\color{mainblue}\Large\bfseries}
  {\thesection}{1em}{}
  [\vspace{-0.5em}\color{mainblue}\rule{\textwidth}{1.2pt}]

\titleformat{\subsection}
  {\color{darkgreen}\large\bfseries}
  {\thesubsection}{0.8em}{}

\titleformat{\subsubsection}
  {\color{darkred}\normalsize\bfseries}
  {\thesubsubsection}{0.6em}{}

% ========== سربرگ ==========
\pagestyle{fancy}
\fancyhf{}
\fancyhead[R]{\small\color{mainblue}حل تمارین اقتصاد کلان پیشرفته}
\fancyhead[L]{\small\color{mainblue}مقطع دکتری}
\fancyfoot[C]{\thepage}
\renewcommand{\headrulewidth}{0.6pt}

% ========== دستورات ریاضی ==========
\newcommand{\R}{\mathbb{R}}
\newcommand{\argmax}{\operatorname*{arg\,max}}
\newcommand{\MRS}{\mathrm{MRS}}
\newcommand{\MRT}{\mathrm{MRT}}

% ========== شروع سند ==========
\begin{document}

\begin{center}
	{\Huge\bfseries\color{mainblue} حل کامل تمارین اقتصاد کلان پیشرفته}\\[0.5cm]
	{\large\color{darkgreen} مقطع دکتری}\\[0.3cm]
	{\color{mainblue}\rule{\textwidth}{1.5pt}}
\end{center}

\tableofcontents
\newpage

% ============================================================
\section{بخش اول — پرسش‌های تحلیلی (۲۵ نمره)}
% ============================================================

% ------------------------------------------------------------
\subsection{سؤال ۱ — قانون والراس (۸ نمره)}
% ------------------------------------------------------------

\subsubsection*{مفهوم کلی مسئله}

تصور کنید در یک اقتصاد بزرگ، چندین کالا و چندین خانوار وجود دارد. هر خانوار مقداری از کالاها را در اختیار دارد که به آن‌ها «موهبت اولیه» می‌گویند. خانوارها می‌خواهند این موهبت‌ها را با کالاهای مورد علاقه‌شان مبادله کنند تا مطلوبیت‌شان بیشینه شود. همزمان، بنگاه‌ها با استفاده از نهاده‌ها کالا تولید می‌کنند و سود می‌برند. سؤال اساسی این است که در چنین اقتصادی، چه ترکیبی از قیمت‌ها و تخصیص کالاها می‌تواند یک «تعادل» ایجاد کند؟ تعادل یعنی وضعیتی که در آن هیچ خانوار و هیچ بنگاهی انگیزه‌ای برای تغییر تصمیم خود نداشته باشد و همه بازارها تسویه شوند.

\subsubsection*{قسمت ۱: تعریف تعادل رقابتی (والراسی)}

تعادل رقابتی یک مفهوم بنیادی در اقتصاد خرد است که توسط لئون والراس معرفی شد. در این تعادل، سه شرط اساسی برقرار است. شرط اول این است که هر خانوار با توجه به قیمت‌های داده‌شده، بهترین سبد مصرفی خود را انتخاب کند. به عبارت دیگر، خانوار مطلوبیت خود را با توجه به محدودیت بودجه‌اش بیشینه می‌کند. بودجه خانوار شامل ارزش موهبت‌های اولیه‌اش به‌علاوه سهمی از سود بنگاه‌هاست که به او تعلق دارد. شرط دوم این است که هر بنگاه با توجه به قیمت‌های داده‌شده، مقدار تولیدی را انتخاب کند که سودش را بیشینه کند. شرط سوم این است که در هر بازار، مجموع تقاضا برابر مجموع عرضه باشد. عرضه شامل موهبت‌های اولیه به‌علاوه تولید بنگاه‌هاست. این سه شرط با هم یک دستگاه معادلات تشکیل می‌دهند که حل آن‌ها تعادل اقتصاد را مشخص می‌کند.

برای خانوار $i$، مسئله بهینه‌یابی به‌صورت زیر است:
\[
	x_i^* \in \argmax_{x_i \geq 0} \ u_i(x_i)
	\quad \text{s.t.} \quad
	p^* \cdot x_i \leq p^* \cdot \omega_i + \sum_{j} \theta_{ij} \, p^* \cdot y_j^*
\]
در این رابطه، $x_i^*$ سبد مصرفی بهینه، $u_i$ تابع مطلوبیت، $p^*$ بردار قیمت‌ها، $\omega_i$ موهبت اولیه، $\theta_{ij}$ سهم خانوار $i$ از بنگاه $j$، و $y_j^*$ تولید بهینه بنگاه $j$ است. برای بنگاه $j$ نیز داریم:
\[
	y_j^* \in \argmax_{y_j \in Y_j} \ p^* \cdot y_j
\]
و شرط تسویه بازار:
\[
	\sum_{i} x_{il}^* = \sum_{i} \omega_{il} + \sum_{j} y_{jl}^*
\]

تفاوت بین «تعریف» و «حل» بسیار مهم است. تعریف به ما می‌گوید که یک تعادل چه ویژگی‌هایی دارد، یعنی چه شرایطی باید برقرار باشد تا بتوانیم بگوییم اقتصاد در تعادل است. اما حل به ما می‌گوید که چگونه می‌توانیم مقادیر عددی قیمت‌ها و تخصیص‌ها را پیدا کنیم. به بیان دیگر، تعریف یک چارچوب مفهومی است و حل یک فرایند عملی. در مسائل اقتصاد کلان، ابتدا باید تعادل را به‌درستی تعریف کنیم و سپس با استفاده از ابزارهای ریاضی آن را حل کنیم.

\subsubsection*{قسمت ۲: قانون والراس}

قانون والراس یکی از مهم‌ترین نتایج در نظریه تعادل عمومی است. این قانون می‌گوید که در یک اقتصاد با ترجیحات محلی غیراشباع، مجموع ارزش پولی مازاد تقاضا در همه بازارها همیشه صفر است. «محلی غیراشباع» یعنی هر خانوار همیشه بتواند با مصرف کمی بیشتر از یک کالا، مطلوبیت خود را افزایش دهد. این فرض بسیار ضعیف و منطقی است، زیرا در دنیای واقعی تقریباً همه افراد همیشه چیزهایی دارند که آرزو می‌کنند بیشتر مصرف کنند.

برای اثبات قانون والراس، ابتدا توجه می‌کنیم که هر خانوار بودجه‌اش را کامل خرج می‌کند. اگر بودجه را کامل خرج نکند، می‌تواند با مصرف کمی بیشتر، مطلوبیت خود را افزایش دهد، که با بهینه بودن سبد مصرفی در تناقض است. بنابراین برای هر خانوار:
\[
	p \cdot x_i = p \cdot \omega_i + \sum_{j} \theta_{ij} \, p \cdot y_j
\]
با جمع روی همه خانوارها:
\[
	\sum_{i} p \cdot x_i = \sum_{i} p \cdot \omega_i + \sum_{i} \sum_{j} \theta_{ij} \, p \cdot y_j
\]
از آنجا که $\sum_{i} \theta_{ij} = 1$ (همه سهام بنگاه‌ها به خانوارها تعلق دارد)، جمله آخر برابر $\sum_j p \cdot y_j$ می‌شود. بنابراین:
\[
	p \cdot \left( \sum_{i} x_i - \sum_{i} \omega_i - \sum_{j} y_j \right) = 0
\]
که همان $p \cdot z(p) = 0$ است. این نتیجه به ما می‌گوید که اگر $L-1$ بازار تسویه شوند، بازار $L$ام خودبه‌خود تسویه می‌شود. پس در حل مسائل، می‌توانیم یکی از معادلات تسویه را حذف کنیم و دستگاه را ساده‌تر حل کنیم. این خاصیت در مسائل تعادل عمومی بسیار کاربردی است.

\subsubsection*{قسمت ۳: یکسان بودن قیمت‌ها}

اگر همه کالاها قیمت یکسانی داشته باشند، شرط انتخاب کالای واحد به‌صورت زیر است:
\[
	\frac{\partial u_i / \partial x_{il}}{\partial u_i / \partial x_{ik}} = 1
\]
این رابطه می‌گوید که نرخ نهایی جانشینی بین هر دو کالا برابر ۱ است. یعنی خانوار حاضر است یک واحد از کالای $l$ را با یک واحد از کالای $k$ مبادله کند. اگر ترجیحات اکیداً محدب نباشند یا موهبت‌های اولیه متفاوت باشند، ممکن است خانوارها کالای واحد را انتخاب نکنند و تخصیص‌های گوشه‌ای رخ دهد. در چنین حالتی، ممکن است قیمت برخی کالاها صفر شود (کالای آزاد). این پدیده با شرایط کوهن-تاکر مطابقت دارد:
\[
	\frac{\partial u_i}{\partial x_{il}} \leq \lambda p_l, \quad x_{il} \geq 0, \quad x_{il} \left( \frac{\partial u_i}{\partial x_{il}} - \lambda p_l \right) = 0
\]
اگر $p_l = 0$ و $\frac{\partial u_i}{\partial x_{il}} > 0$، آنگاه $x_{il} > 0$ و کالا مصرف می‌شود.

% ------------------------------------------------------------
\subsection{سؤال ۲ — قضایای رفاه (۹ نمره)}
% ------------------------------------------------------------

\subsubsection*{مفهوم کلی مسئله}

قضایای رفاه به ما می‌گویند که آیا بازار رقابتی نتایج «خوب» تولید می‌کند یا نه. «خوب» در اینجا به معنای «پارتوکارا» است، یعنی تخصیصی که نمی‌توان کسی را بهتر کرد بدون اینکه کس دیگری بدتر شود. این قضایا پایه‌های نظری اقتصاد رفاه را تشکیل می‌دهند و به سیاست‌گذاران کمک می‌کنند تا بدانند چه زمانی بازار کارا عمل می‌کند و چه زمانی نیاز به دخالت دارد.

\subsubsection*{قسمت ۱: قضیه اول و دوم رفاه}

قضیه اول رفاه می‌گوید که هر تعادل رقابتی، پارتوکارا است. این قضیه بسیار قوی است زیرا نشان می‌دهد که بازار رقابتی به‌طور خودکار به کارایی می‌رسد، بدون نیاز به دخالت دولت. فرض‌های این قضیه شامل ترجیحات محلی غیراشباع، رقابت کامل، اطلاعات کامل، عدم وجود آثار خارجی و عدم وجود کالاهای عمومی است. اثبات این قضیه با استفاده از روش برهان خلف انجام می‌شود. فرض کنید تخصیص تعادلی $x^*$ پارتوکارا نباشد. آنگاه تخصیص دیگری $x'$ وجود دارد که برای همه خانوارها حداقل به‌اندازه $x^*$ مطلوب است و برای برخی اکیداً بهتر است. چون $x_i^*$ در بودجه $p^*$ بهینه است، $x_i'$ باید گران‌تر از موهبت خانوار $i$ باشد. با جمع روی $i$ و استفاده از تسویه بازار، به تناقض می‌رسیم.

قضیه دوم رفاه می‌گوید که هر تخصیص پارتوکارا را می‌توان به‌عنوان یک تعادل رقابتی با بازتوزیع موهبت‌های اولیه به‌دست آورد. این قضیه اهمیت بسیاری دارد زیرا نشان می‌دهد که می‌توان با انتقال یکجای ثروت، توزیع را تغییر داد بدون اینکه کارایی از بین برود. فرض‌های این قضیه شامل تحدب ترجیحات، تحدب مجموعه‌های تولیدی، وجود موهبت‌های اولیه قابل انتقال و پیوستگی ترجیحات است. تحدب برای این قضیه ضروری است زیرا باید بتوانیم یک ابرصفحه جداکننده بین مجموعه‌های محدب پیدا کنیم تا قیمت‌ها را استخراج کنیم. بدون تحدب، ممکن است تخصیص پارتوکارا با قیمت‌های خطی پشتیبانی نشود.

تفاوت دو قضیه را می‌توان در جدول زیر خلاصه کرد:

\begin{center}
	\begin{tabular}{|c|c|c|}
		\hline
		\textbf{ویژگی} & \textbf{قضیه اول}    & \textbf{قضیه دوم}        \\
		\hline
		جهت            & از تعادل به کارایی   & از کارایی به تعادل       \\
		\hline
		نقش بازتوزیع   & ندارد                & دارد (انتقال یکجا)       \\
		\hline
		فرض تحدب       & لازم نیست            & ضروری است                \\
		\hline
		کاربرد         & ارزیابی کارایی بازار & طراحی سیاست‌های بازتوزیعی \\
		\hline
	\end{tabular}
\end{center}

\subsubsection*{قسمت ۲: تعادل رقابتی و عدالت}

اگر تخصیص $x^*$ یک تعادل رقابتی با بازتوزیع موهبت‌ها باشد و ترجیحات یکسان و محدب باشند، آنگاه $x^*$ عادلانه است. عدالت در اینجا به معنای این است که هیچ‌کس سبد دیگری را ترجیح نمی‌دهد. به بیان دیگر، همه از سبد خودشان راضی هستند و انگیزه‌ای برای مبادله ندارند. اثبات این گزاره ساده است: فرض کنید کسی ترجیح دهد سبد دیگری $x_i'$ را که از نظر ارزش برابر است، داشته باشد. آنگاه با بهینه‌یابی خانوار در تناقض است، چون $x_i^*$ باید حداقل به‌اندازه $x_i'$ مطلوبیت بدهد. بنابراین هیچ‌کس سبد دیگری را ترجیح نمی‌دهد. با انتقال یکجا موهبت‌ها، تعادل جدید رقابتی است و تخصیص پارتوکارای مورد نظر را پشتیبانی می‌کند.

\subsubsection*{قسمت ۳: سه حالت مقایسه}

سه حالت مختلف را می‌توان مقایسه کرد. حالت اول انتقال یکجاست که در آن قیمت‌ها تغییر نمی‌کنند و اعوجاجی ایجاد نمی‌شود. نرخ نهایی جانشینی با نرخ نهایی تبدیل برابر می‌ماند. حالت دوم مالیات نسبی بر دستمزد است که نرخ نهایی جانشینی بین کار و فراغت را از نرخ نهایی تبدیل فنی جدا می‌کند:
\[
	\MRS_{l,c} = \frac{w(1 - \tau)}{p} \neq \frac{w}{p} = \MRT_{l,c}
\]
این شکاف، بار اضافی مالیات یا Deadweight Loss ایجاد می‌کند. حالت سوم اثر خارجی تولید است که در آن اگر بنگاه هزینه اجتماعی تولید را نپردازد، قیمت خصوصی با هزینه اجتماعی برابر نیست:
\[
	p \neq MSC = MPC + MEC
\]
و تخصیص بازار کارا نخواهد بود. در جدول زیر این سه حالت مقایسه شده‌اند:

\begin{center}
	\begin{tabular}{|c|c|c|c|}
		\hline
		\textbf{حالت} & \textbf{کارایی} & \textbf{توزیع} & \textbf{قضیه اول} \\
		\hline
		انتقال یکجا   & حفظ می‌شود       & مستقیم         & برقرار است        \\
		\hline
		مالیات نسبی   & اعوجاج          & غیرمستقیم      & نقض می‌شود         \\
		\hline
		اثر خارجی     & شکست بازار      & بستگی دارد     & نقض می‌شود         \\
		\hline
	\end{tabular}
\end{center}

% ------------------------------------------------------------
\subsection{سؤال ۳ — معادله اویلر (۸ نمره)}
% ------------------------------------------------------------

\subsubsection*{مفهوم کلی مسئله}

خانوار می‌خواهد بین مصرف امروز و مصرف فردا تعادل برقرار کند. اگر امروز کمتر مصرف کند، می‌تواند پس‌انداز کند و فردا بیشتر مصرف کند. معادله اویلر این تصمیم را توصیف می‌کند و یکی از مهم‌ترین معادلات در اقتصاد کلان است. این معادله می‌گوید که مطلوبیت نهایی مصرف امروز باید برابر باشد با مطلوبیت نهایی مصرف فردا که با عامل تخفیف $\beta$ و نرخ بهره $R$ تعدیل شده است.

\subsubsection*{قسمت ۱: رابطه مصرف بین دو دوره}

مسئله خانوار به‌صورت زیر است:
\[
	\max_{c_1, c_2} \ u(c_1) + \beta u(c_2)
	\quad \text{s.t.} \quad
	c_1 + \frac{c_2}{R} = W
\]
خانوار می‌خواهد مطلوبیت کل را بیشینه کند. مطلوبیت کل شامل مطلوبیت مصرف امروز به‌علاوه مطلوبیت مصرف فردا با تخفیف $\beta$ است. قید بودجه می‌گوید که ارزش حال مصرف امروز و فردا برابر ثروت کل $W$ است. شرط مرتبه اول به‌صورت زیر است:
\[
	u'(c_1) = \beta R \, u'(c_2)
\]
این رابطه می‌گوید که مطلوبیت نهایی یک واحد مصرف امروز باید برابر باشد با مطلوبیت نهایی یک واحد مصرف فردا که با $\beta R$ تعدیل شده است. با تابع مطلوبیت $u(c) = \frac{c^{1-\sigma}}{1-\sigma}$، داریم:
\[
	c_1^{-\sigma} = \beta R \, c_2^{-\sigma}
	\quad \Rightarrow \quad
	\frac{c_2}{c_1} = (\beta R)^{1/\sigma}
\]
پارامتر $\sigma$ کشش جانشینی بین‌دوره‌ای معکوس است. اگر $\sigma$ کوچک باشد، کشش جانشینی بزرگ است و مصرف نسبت به تغییرات $R$ حساس‌تر است. اگر $\sigma = 1$ (لگاریتمی)، $\frac{c_2}{c_1} = \beta R$. اگر $\sigma \to \infty$ (لئونتیف)، مصرف بین دو دوره ثابت می‌ماند.

\subsubsection*{قسمت ۲: اثر افزایش نرخ بهره}

افزایش نرخ بهره $R$ دو اثر متقابل دارد. اثر جانشینی می‌گوید که مصرف آینده نسبت به حال ارزان‌تر می‌شود، پس مصرف دوره ۱ کاهش و پس‌انداز افزایش می‌یابد. اثر ثروت می‌گوید که اگر خانوار وام‌دهنده خالص باشد، افزایش $R$ ثروت او را افزایش می‌دهد و مصرف هر دو دوره را بالا می‌برد. اگر وام‌گیرنده خالص باشد، اثر ثروت منفی است. اگر $\sigma < 1$، اثر جانشینی غالب است. اگر $\sigma > 1$، اثر ثروت ممکن است غالب شود. اگر $\sigma = 1$، اثرات متعادل می‌شوند و $\frac{\partial c_1}{\partial R} = 0$.

% ============================================================
\section{بخش دوم — مسائل (۷۵ نمره)}
% ============================================================

% ------------------------------------------------------------
\subsection{سؤال ۴ — اقتصاد مبادله‌ای (۲۰ نمره)}
% ------------------------------------------------------------

\subsubsection*{مفهوم کلی مسئله}

دو خانوار داریم که دو کالا (x و y) مصرف می‌کنند. هر خانوار مقداری از این کالاها را در اختیار دارد و می‌خواهد با مبادله، مطلوبیت خود را بیشینه کند. سؤال این است که قیمت تعادلی و تخصیص تعادلی چیست. این مسئله یک مثال کلاسیک از تعادل عمومی در اقتصاد مبادله‌ای است.

\subsubsection*{قسمت ۱: تقاضای والراسی}

ترجیحات دو خانوار به‌صورت کاب-داگلاس است:
\[
	u_i(x_i, y_i) = \alpha_i \ln x_i + (1 - \alpha_i) \ln y_i
\]
موهبت‌ها و پارامترها به‌صورت زیر است:
\[
	\omega_1 = (4, 1), \quad \omega_2 = (1, 3)
\]
\[
	\alpha_1 = \frac{3}{4}, \quad \alpha_2 = \frac{1}{4}
\]
در ترجیحات کاب-داگلاس، خانوار بخش ثابتی از درآمدش را به هر کالا اختصاص می‌دهد. مثلاً خانوار ۱ که $\alpha_1 = 3/4$ دارد، ۷۵٪ درآمدش را به کالای x می‌دهد. تقاضای مارشالی برای هر کالا به‌صورت زیر است:
\[
	x_i = \alpha_i (\omega_{xi} + q \omega_{yi}), \quad
	y_i = \frac{1 - \alpha_i}{q} (\omega_{xi} + q \omega_{yi})
\]
برای خانوار ۱:
\[
	x_1 = \frac{3}{4}(4 + q), \quad y_1 = \frac{1}{4q}(4 + q)
\]
برای خانوار ۲:
\[
	x_2 = \frac{1}{4}(1 + 3q), \quad y_2 = \frac{3}{4q}(1 + 3q)
\]

\subsubsection*{قسمت ۲: تعادل رقابتی}

تسویه بازار x به‌صورت زیر است:
\[
	x_1 + x_2 = 5
\]
\[
	\frac{3}{4}(4 + q) + \frac{1}{4}(1 + 3q) = 5
\]
\[
	3 + \frac{3q}{4} + \frac{1}{4} + \frac{3q}{4} = 5
\]
\[
	\frac{13}{4} + \frac{3q}{2} = 5
	\quad \Rightarrow \quad
	q^* = \frac{7}{6}
\]
تخصیص تعادلی به‌صورت زیر است:
\[
	x_1 = \frac{31}{8}, \quad x_2 = \frac{9}{8}, \quad y_1 = \frac{31}{28}, \quad y_2 = \frac{81}{28}
\]
کنترل تسویه بازار y:
\[
	y_1 + y_2 = \frac{31}{28} + \frac{81}{28} = 4 = \omega_{y1} + \omega_{y2}
\]

\subsubsection*{قسمت ۳: بهینگی پارتو}

شرط کارایی در اقتصاد مبادله‌ای این است که نرخ نهایی جانشینی دو خانوار برابر باشد:
\[
	\MRS^1 = \frac{3y_1}{x_1}, \quad \MRS^2 = \frac{y_2}{3x_2}
\]
در تعادل رقابتی، $\MRS = \frac{1}{q} = \frac{6}{7}$. بررسی:
\[
	\frac{3 \cdot 31/28}{31/8} = \frac{6}{7}, \quad
	\frac{81/28}{3 \cdot 9/8} = \frac{6}{7}
\]
پس تخصیص تعادلی پارتوکارا است و قضیه اول رفاه تأیید می‌شود.

\subsubsection*{قسمت ۴: انتقال یکجای ثروت}

اگر $t$ واحد از خانوار ۱ به خانوار ۲ منتقل شود، موهبت‌های جدید به‌صورت زیر می‌شوند:
\[
	\omega_1' = (4 - t, 1), \quad \omega_2' = (1 + t, 3)
\]
با تسویه بازار x:
\[
	q(t) = \frac{7}{6} + \frac{t}{3}
\]
این نتیجه نشان می‌دهد که با انتقال ثروت، قیمت نسبی تغییر می‌کند چون تقاضای نسبی دو خانوار متفاوت است. در $t = 0$، قیمت به تعادل اولیه بازمی‌گردد.

\subsubsection*{قسمت ۵: استقلال قیمت از توزیع}

اگر $\alpha_1 = \alpha_2 = \alpha$، آنگاه:
\[
	q = \frac{\Omega_x (1 - \alpha)}{\alpha \Omega_y}
\]
قیمت نسبی فقط به ترجیحات و موهبت کل بستگی دارد، نه به توزیع ثروت. این نتیجه‌ای کلاسیک است: با ترجیحات یکسان و همگن، قیمت‌های نسبی از توزیع مستقل می‌شوند.

% ------------------------------------------------------------
\subsection{سؤال ۵ — اقتصاد تولیدی (۲۰ نمره)}
% ------------------------------------------------------------

\subsubsection*{مفهوم کلی مسئله}

یک اقتصاد یک‌دوره‌ای داریم. خانوارها کار می‌کنند و دستمزد می‌گیرند، بنگاه‌ها با کار تولید می‌کنند و سود می‌دهند. سؤال این است که تعادل کار و تولید چیست. این مسئله به ما نشان می‌دهد که چگونه عرضه کار و تقاضای کار در یک اقتصاد ساده با هم تعادل برقرار می‌کنند.

\subsubsection*{قسمت ۱: حل یکتا}

مسئله خانوار به‌صورت زیر است:
\[
	\max \ \ln c + \chi \ln(1-h)
	\quad \text{s.t.} \quad
	c = w h + \pi
\]
خانوار مطلوبیت مصرف و فراغت را بیشینه می‌کند. قید بودجه می‌گوید که مصرف برابر دستمزد ضربدر کار به‌علاوه سود بنگاه است. مسئله بنگاه:
\[
	\max_{H} \ A H^\theta - w H
	\quad \Rightarrow \quad
	w = A \theta H^{\theta - 1}
\]
بنگاه سود خود را بیشینه می‌کند. شرط مرتبه اول تقاضای کار را می‌دهد. با حل همزمان این دو مسئله، تعادل به‌صورت زیر است:
\[
	H = \frac{\theta}{\theta + \chi}
\]
عرضه کار فقط به $\theta$ (سهم کار در تولید) و $\chi$ (ارزش فراغت) بستگی دارد.

\subsubsection*{قسمت ۲: تخصیص تعادل}

تخصیص تعادل به‌صورت زیر است:
\[
	w = A \theta \left(\frac{\theta}{\theta + \chi}\right)^{\theta - 1}, \quad
	c = A \left(\frac{\theta}{\theta + \chi}\right)^\theta, \quad
	1 - H = \frac{\chi}{\theta + \chi}
\]
این روابط نشان می‌دهند که چگونه دستمزد، مصرف و فراغت در تعادل تعیین می‌شوند.

\subsubsection*{قسمت ۳: اثر $A$ و $\chi$}

\begin{center}
	\begin{tabular}{|c|c|c|c|}
		\hline
		\textbf{پارامتر} & $H$   & $w$    & $c$    \\
		\hline
		افزایش $A$       & بی‌اثر & افزایش & افزایش \\
		\hline
		افزایش $\chi$    & کاهش  & افزایش & کاهش   \\
		\hline
	\end{tabular}
\end{center}

افزایش $A$ (بهره‌وری) چون ترجیحات لگاریتمی است، اثر جانشینی و ثروت خنثی می‌شوند و $H$ تغییر نمی‌کند. اما $w$ و $c$ افزایش می‌یابند. افزایش $\chi$ (ارزش فراغت) باعث می‌شود مردم کمتر کار کنند، دستمزد افزایش یابد و مصرف کاهش یابد.

\subsubsection*{قسمت ۴: مالیات $\tau$}

با مالیات $\tau$ بر درآمد و بازپرداخت $T = \tau w H$:
\[
	H = \frac{\theta (1 - \tau) + \tau \theta}{\theta (1 + \chi) - \tau \theta \chi + \chi \theta}
\]
مالیات بر درآمد اعوجاج ایجاد می‌کند و عرضه کار را کاهش می‌دهد. انتقال یکجا این اعوجاج را ندارد.

\subsubsection*{قسمت ۵: مقایسه مالیات‌ها}

مالیات بر سود با بازپرداخت یکجا اعوجاج کمتری دارد چون دستمزد (قیمت نسبی) را تغییر نمی‌دهد. در مقابل، مالیات دستمزد نرخ جانشینی بین مصرف و فراغت را تغییر می‌دهد و اعوجاج ایجاد می‌کند.

% ------------------------------------------------------------
\subsection{سؤال ۶ — اقتصاد دودوره‌ای (۲۰ نمره)}
% ------------------------------------------------------------

\subsubsection*{مفهوم کلی مسئله}

دو خانوار A و B داریم. A در دوره ۱ درآمد زیاد (۳) و در دوره ۲ درآمد کم (۱) دارد. B برعکس. آن‌ها می‌توانند وام بگیرند یا وام بدهند تا مصرف خود را هموار کنند. سؤال این است که نرخ بهره تعادلی و تخصیص تعادلی چیست. این مسئله به ما نشان می‌دهد که چگونه بازارهای مالی به هموارسازی مصرف کمک می‌کنند.

\subsubsection*{قسمت ۱: تقاضای دارایی}

قید بودجه هر خانوار به‌صورت زیر است:
\[
	c_{i1} + a_i = y_{i1}, \quad c_{i2} = y_{i2} + (1 + r) a_i
\]
خانوار در دوره ۱ مقداری پس‌انداز ($a_i$) می‌کند و در دوره ۲ آن را با بهره دریافت می‌کند. بهینه‌یابی:
\[
	a_i = \frac{y_{i1}}{2} - \frac{y_{i2}}{2(1 + r)}
\]
برای خانوار A: $a_A = \frac{3}{2} - \frac{1}{2(1+r)}$
برای خانوار B: $a_B = \frac{1}{2} - \frac{3}{2(1+r)}$

\subsubsection*{قسمت ۲: بدون محدودیت}

تسویه بازار دارایی $a_A + a_B = 0$:
\[
	r = 0, \quad c_{A1} = c_{A2} = c_{B1} = c_{B2} = 2
\]
چون A و B می‌توانند آزادانه وام بگیرند/بدهند، مصرف خود را کاملاً هموار می‌کنند.

\subsubsection*{قسمت ۳: محدودیت $a_B \geq 0$}

B نمی‌تواند قرض بگیرد. پس $a_B = 0$:
\[
	c_{B1} = 1, \quad c_{B2} = 3, \quad c_{A1} = 3, \quad c_{A2} = 1
\]
\[
	r = -\frac{2}{3}
\]
محدودیت استقراض فعال است و تعادل جدید ایجاد می‌شود.

\subsubsection*{قسمت ۴: ناپارتوکارا}

$\MRS_A = 1/3$ و $\MRS_B = 3$. چون برابر نیستند، تخصیص پارتوکارا نیست. بهبود پارتو: انتقال ۱ واحد از A به B در دوره ۱ و برعکس در دوره ۲:
\[
	c_A' = (2, 2), \quad c_B' = (2, 2)
\]
هر دو بهتر می‌شوند.

\subsubsection*{قسمت ۵: هر دو ممنوع}

اگر هر دو ممنوع از استقراض باشند:
\[
	-\frac{2}{3} \leq r \leq 2
\]
هر نرخ بهره در این بازه، تعادل است چون هیچ‌کس مایل به معامله نیست.

% ------------------------------------------------------------
\subsection{سؤال ۷ — افق نامحدود (۱۵ نمره)}
% ------------------------------------------------------------

\subsubsection*{مفهوم کلی مسئله}

خانوار در زمان نامحدود زندگی می‌کند. در هر دوره درآمد $y_t$ می‌گیرد و می‌تواند پس‌انداز کند. سؤال این است که مسیر بهینه مصرف چیست. این مسئله یکی از مسائل کلاسیک در اقتصاد کلان است که به ما نشان می‌دهد چگونه خانوارها مصرف خود را در طول زمان هموار می‌کنند.

\subsubsection*{قسمت ۱: محدودیت ارزش حال}

با جمع‌زدن بودجه‌ها:
\[
	\sum_{t=0}^{\infty} \frac{c_t}{R^t} \leq W = R b_0 + \sum_{t=0}^{\infty} \frac{y_t}{R^t}
\]
شرط غیرپانزی $\lim_{T \to \infty} \frac{b_{T+1}}{R^T} \geq 0$ تضمین می‌کند که خانوار در نهایت بدهی‌اش را پرداخت می‌کند و از طرح‌های پانزی جلوگیری می‌شود.

\subsubsection*{قسمت ۲: معادله اویلر و مسیر مصرف}

معادله اویلر:
\[
	\frac{c_{t+1}}{c_t} = \beta R
\]
مسیر مصرف:
\[
	c_t = (1 - \beta) (\beta R)^t W
\]
این مسیر نشان می‌دهد که مصرف در هر دوره به ثروت کل و پارامترهای $\beta$ و $R$ بستگی دارد.

\subsubsection*{قسمت ۳: مقایسه افزایش موقت و دائمی}

افزایش موقت فقط $y_t$ را افزایش می‌دهد. چون درآمد در یک دوره است، اثرش کوچک است. افزایش دائمی همه $y_t$ را افزایش می‌دهد. اثرش بزرگ‌تر است چون ثروت دائمی را افزایش می‌دهد.

\subsubsection*{قسمت ۴: ناکافی بودن اویلر}

معادله اویلر فقط نسبت مصرف دو دوره را می‌دهد، نه سطح آن. برای تعیین سطح، به شرط تراگذری و محدودیت بودجه نیاز داریم. مثلاً مسیر $c_t = C(\beta R)^t$ برای هر $C > 0$ معادله اویلر را برآورده می‌کند، اما فقط یک $C$ محدودیت بودجه را برآورده می‌کند.

\end{document}
